\section{New Student Wording: Design and Annotation}
\label{app:cue-transfer}

\begin{table}[t]
\centering\small
\begin{tabular}{llrr}
\toprule
Expression & Prediction comparison & Gain [95\% CI] & Holm $p$\\
\midrule
Explicit & Default: reveal & .09730 [.08095, .11320] & $7.63\!\times\!10^{-11}$\\
Explicit & Default: elicit & .13070 [.12266, .13825] & $4.97\!\times\!10^{-24}$\\
Explicit & State: acknowledge & -.00366 [-.00793, .00090] & 1.0\\
Implicit & Default: reveal & .08231 [.06657, .09695] & $2.76\!\times\!10^{-9}$\\
Implicit & Default: elicit & .12748 [.11660, .13806] & $8.51\!\times\!10^{-20}$\\
Implicit & State: acknowledge & -.00628 [-.00960, -.00273] & 1.0\\
\bottomrule
\end{tabular}
\caption{Six fixed-prediction primary expression-transfer tests. Gains are
source-averaged reductions in Brier loss; intervals bootstrap 32 reused source
problems, conditioning on fitted predictions. The default baseline is educational
context; the state baseline is the model-by-subject profile. One-sided tests
ask whether the gain is positive: $p=1$ is not a separately specified test of
harm, and the confidence intervals (CIs) are not simultaneous intervals.}
\label{tab:cue-primary}
\end{table}

\begin{table}[t]
\centering\small
\begin{tabular}{p{.12\linewidth}p{.40\linewidth}p{.40\linewidth}}
\toprule
Condition / pair & Calm cue & Frustrated cue\\
\midrule
Original & I feel calm while working on this. & I feel frustrated while working on this.\\
Explicit 1 & I'm feeling at ease as I work through this. & I'm feeling annoyed as I work through this.\\
Explicit 2 & Working on this leaves me feeling relaxed. & Working on this leaves me feeling irritated.\\
Explicit 3 & I'm feeling composed about working on this problem. & I'm feeling exasperated about working on this problem.\\
Explicit 4 & I feel settled while I work on this. & I feel fed up while I work on this.\\
Implicit 1 & I've set aside some time for this; one step at a time is fine. & Ugh, this again. I'm about ready to close the page.\\
Implicit 2 & I'm taking my time with the page. There's no rush. & I've just pushed the page away. Ugh, enough already.\\
Implicit 3 & I'm sitting back with my notebook open, ready to take the next step. & I've thrown my pencil down beside the notebook. Ugh.\\
Implicit 4 & I'll take a slow breath and carry on with the page in front of me. & I'm gritting my teeth at the page in front of me. Ugh.\\
Peer & I feel calm while working on this. My classmate, who is doing a different problem, says, ``I feel calm while working on this.'' & I feel calm while working on this. My classmate, who is doing a different problem, says, ``I feel frustrated while working on this.''\\
\bottomrule
\end{tabular}
\caption{Complete fixed expression stimuli. Calm and frustrated label the intended input conditions, not diagnosed
human emotions. Within each domain, sources
are ordered by their frozen hash: explicit pair index cycles through 1--4;
implicit pair index is offset by one modulo four. Each source keeps its assigned
pair across student-work states, models, and requests.}
\label{tab:exact-cues}
\end{table}

The extension protocol and prediction table were frozen on 7 September 2026
before its first tutor generation. Its 32 problem sources had already been
observed in the first test stage; they are not a second independent
source holdout. Each of the four wording conditions combines 32 problems, two student-work
states, calm/frustrated cues, five models, and two requests, yielding 1,280 answers.
Explicit and implicit conditions each assign one of four fixed phrase pairs per source,
balanced with two sources per domain and phrase pair. Student-work text and
reference answers remain unchanged. The classmate condition assigns the current student the
calm state in the fixed predictor regardless of the peer's quote. Exact
materials, the 61,440-row prediction table covering all outcomes and predictors,
and their hashes are retained with the executable analysis.

The extension yields all 5,120 answers and uses the original three event annotators
and evidence-line rubric. After the recovery described below, all 15,360 codes
are usable. Pairwise positive annotator agreement is 99.25--99.41\% for revelation,
90.68--92.34\% for elicitation, and 94.89--95.51\% for acknowledgement.
Of 15,360 requested codes, 15,359 pass the original strict parser after the
fixed repair budget. The remaining final response contains all eight complete
event objects but lacks one root closing brace. A separately recorded amendment
appends exactly that brace to a derived copy, preserving raw responses, earlier
failed attempts, and the failed controller's exit record. No event label or
evidence is edited and no additional API request is made. The other two annotators
agree on all three primary events for this response; replacing the missing
vote by either binary value leaves every primary majority unchanged. Two
secondary events do not have this invariance. Derived measurement selection,
original diagnostics, and recovered diagnostics are stored separately.

Generation completed on 7 September and measurement recovery on 8 September
2026. The first-stage fitted predictors remain unchanged. Six new-expression
primary tests form their own Holm family; same-period original-wording
comparisons, new-minus-original wording paired contrasts, phrase subgroups,
classmate contrasts, repetition, and checks using alternative annotations or excluding models are secondary.
Bootstrap intervals use 10,000 source resamples stratified by domain with seed
20260907. They condition on the fixed training fit and the finite source/phrase
assignment. The annotator checks use each annotator's labels separately. The model
checks exclude each generating model's responses without refitting the
original predictors.

\paragraph{Original wording rerun during the same period.}
The acknowledgement state gain is +0.01248 with a 95\% source-bootstrap
interval of [0.00888, 0.01624] when models generate new responses to the original wording.
Paired differences in gain between new wording and the original wording rerun are -0.01614
[-0.02254, -0.00990] for explicit synonyms and -0.01876
[-0.02384, -0.01358] for implicit cues. These secondary contrasts compare
expression conditions during the same observation period; they do not isolate
the effects of provider changes over time.

\paragraph{Classmate control.}
Keeping the current student explicitly calm, changing a peer's quote yields
acknowledgement-rate changes of +30.47, +25.00, +26.56, +23.44, and -4.69
percentage points for M2.7, M3, DeepSeek, GLM, and Doubao, respectively. The
first four source intervals exclude zero; Doubao's includes it. These are
secondary descriptive contrasts. For each model, a fixed selection rule chooses one example that changes
from no acknowledgement to acknowledgement: the first four models explicitly affirm
the current student's calmness, sometimes correctly distinguishing the peer's
frustration. The event-rate contrast therefore cannot itself establish
misattribution. Outcome-selected examples do not estimate its frequency.


\section{Reconstructing Historical Inputs and Their Limits}

A local reconstruction of six tutor-dialogue settings recovers 3,454 task items
and 960 exact question/conversation source groups, including 940 groups shared
across benchmark settings. Of these, 551 items in 250 groups lack a separately
recovered target question; the remaining 710 groups have at least one known
problem. Reference fields comprise 2,679 provided non-placeholder solutions,
475 placeholders, and 300 absent references. Provided solutions are not thereby
verified as correct. Exact grouping normalizes Unicode and whitespace, not
mathematical content or paraphrase equivalence, and does not certify statistical
independence.

All 92 prediction-file and summary pairs match the earlier census snapshots.
The reconstruction links 54,264 latest-success file-item records without
unmatched joins, yielding 46,356 distinct variant-preserving response-text
records. The standard task subset for seven models contains 24,178 records.
Reviewed historical adapter methods support reconstruction consistency, but do
not establish all historic data, constants, client behavior, or deployed weights.

A separate protocol selects 256 eligible exact source groups with paired
scaffolding/pedagogy replies from seven historical model aliases: 3,584 existing
answers, requiring 10,752 automated codes. Source-grouped cross-fitting and
paired instruction effects were specified before coding. Coding on 8 September
2026 produced 10,750 valid codes; two provider-blocked codes remain absent.
The requirement for complete coding failed, but the two available votes identify every
affected majority regardless of the missing vote. Appendix~\ref{app:archive-validation}
reports the resulting analysis and its post-hoc missing-data amendment.

An input-only scope audit reproduces the frozen selection exactly. The standard
scaffolding pool has 1,150 items in 741 exact source groups; excluding 148 items
with missing problems and placeholder references leaves 1,002 items in 612
groups. Source-level selection chooses 256 groups and one dialogue per group.
The selected contexts contain 243 two-utterance, twelve three-utterance, and one
four-utterance histories. Median context lengths are 398 characters in the full
pool, 422 in the eligible pool, and 429 in the selection. These summaries do not
establish representativeness of classroom interactions or the full archive.

MathTutorBench combines Bridge's real tutoring snippets
\citep{wang2024bridge} and MathDial's teacher--simulated-student dialogues
\citep{macina2023mathdial} into its MathDialBridge collection
\citep{macina2025mathtutorbench}. All selected problem strings exactly match
questions in the local MathDial inventory, but bridge records lack per-item
upstream dataset identifiers. Matching questions do not establish where each dialogue originated or whether
a real learner participated. Accordingly, we refer to archived
benchmark dialogue rather than real-classroom validation. Existing inputs and
references may be human-authored; upstream human quality/preference labels and
the omitted target tutor utterance are not used as behavioral ground truth.



\section{Explaining the Prediction Changes After Inspecting Results}
\label{app:cue-diagnosis}

This diagnostic was specified after the extension's primary results and
absolute-loss summaries had been inspected. It adds no API calls, fitted
parameters, hypothesis tests, or independent test. The three student
wording conditions share identical fixed predictions at every matched source,
domain, model, student-work state, calm or frustrated cue, and repetition index. Classmate controls
are excluded because their own-student affect assignments differ.

For baseline probability $b$, conditional probability $e$, contemporary label
$y_0$, and new-expression label $y_1$, let $d=e-b$ and $z=y_1-y_0$.
The change in Brier-loss gain has the identity for the observed data
\begin{equation}
\Delta G = 2\operatorname{E}[zd]
         = 2\operatorname{E}[z]\operatorname{E}[d] + 2\operatorname{Cov}(z,d).
\end{equation}
Expectations and covariance use equally weighted observed cells, not a causal
model. For acknowledgement, mean conditional-minus-domain probability is
0.000009812. The overall-rate terms are approximately -0.000001043 for explicit
synonyms and -0.000002162 for implicit cues, compared with total gain changes
of -0.016144 and -0.018761. A common frequency shift alone therefore cannot
account for this relative contrast. This does not evaluate the generalization
of any recalibrated predictor or exclude model/cue-specific recalibration.

\begin{table}[ht]
\centering\small
\begin{tabular}{lrrrr}
\toprule
Model & Frozen cue gap & Original & Explicit & Implicit\\
\midrule
MiniMax-M2.7 & .5417 & .3516 & .4375 & .4531\\
MiniMax-M3 & .3597 & .3281 & .7656 & .7500\\
DeepSeek-V4-Pro & .2625 & .1875 & .4844 & .5078\\
Doubao-Seed-2.0-Lite & .7455 & .7500 & .7422 & .6484\\
GLM-5.3 & .6788 & .6094 & .6406 & .8438\\
\bottomrule
\end{tabular}
\caption{Acknowledgement-rate gaps: frustrated condition minus calm condition, averaged
across matched problems, work states, and requests. The fixed conditional
probability gap is identical across wording conditions; the other columns are
observed gaps. These are post-result descriptive contrasts without new tests.}
\end{table}

M3 and DeepSeek show larger, not absent, acknowledgement cue gaps under the new
expressions. Their weighted contributions to the total explicit gain change
are -0.005577 and -0.010557, and to the implicit change -0.005382 and -0.013470.
M2.7 and GLM contribute in the positive direction, and Doubao in the negative
direction. Full model, domain, work-state, and calm/frustrated cue contributions are
retained and sum to the respective total within numerical precision. These
are descriptions of the same outputs, not causal effects or independent
replications. New and contemporary replies need not share random seeds for
the accounting identity; their pairing does not represent a real learner's
trajectory. The two wording conditions' different absolute-loss changes and the
finite cue bank continue to limit a single-mechanism interpretation.


\section{Testing Recalibration on Held-Out Problems}
\label{app:cue-calibration}

After inspecting the primary results and absolute-loss diagnostic, we specify a
separate acknowledgement-only check. The domain and conditional predictors
receive identical calibration procedures. The identity map retains original
probabilities. Shared logistic calibration fits $\sigma(a+b\operatorname{logit}(p))$
across models; model-specific calibration fits one such map per model.
We clip input probabilities to $[10^{-6},1-10^{-6}]$ for logistic maps, minimize
the summed binomial negative log likelihood plus
$\tfrac12[a^2+(b-1)^2]$, and fix $a\in[-10,10]$, $b\in[0,5]$ without tuning.
All 240 fitted maps converge without reaching these bounds.

Each domain's eight source families are hash-ordered and assigned cyclically
to four folds. Every model, state, repetition, and expression for a source
shares its fold. Each fit uses 24 sources and predicts eight held-out sources.
One calibration setting trains only on new responses to the original wording;
the other trains on responses to the wording being evaluated. The
latter evaluates adaptation to source-held-out inputs, not zero-shot expression
transfer. Training and test problems are separate, but may use the same phrases. Tests verify that held-out source labels cannot change their
predictions and target-condition labels cannot affect original-wording calibration.

\begin{table}[ht]
\centering\small
\begin{tabular}{llrrr}
\toprule
Training condition & Map & Original & Explicit & Implicit\\
\midrule
None & Identity & .012483 & -.003661 & -.006277\\
Original & Shared logistic & .012130 & -.002864 & -.005871\\
Original & Per-model logistic & .005646 & -.000846 & -.003864\\
Target & Shared logistic & --- & -.005239 & -.010456\\
Target & Per-model logistic & --- & -.003534 & -.008231\\
\bottomrule
\end{tabular}
\caption{Post-result acknowledgement Brier gains: calibrated domain loss minus
calibrated conditional loss, with the same map class applied to both. Original
means original wording rerun during the same period. Target-condition calibration uses target labels
from other sources. Negative estimates do not all have intervals excluding zero.}
\end{table}

Under target-condition per-model calibration, explicit conditional loss improves from
0.144501 to 0.132657, while domain loss improves from 0.140840 to 0.129123.
Implicit conditional loss improves from 0.155411 to 0.133293, while domain loss
improves from 0.149134 to 0.125062. Adaptation can therefore improve absolute
prediction without restoring the advantage of modeling each model's
response to student cues.
Per-model calibration can itself absorb model-dependent cue responses, so
successful baseline adaptation is not evidence of their absence.

All absolute losses, parameters, source predictions, and domain-stratified
source bootstrap intervals are retained. The 10,000 draws use seed 20260908;
intervals condition on cross-fitted predictions without refitting each bootstrap
or accounting for the full post-result selection history. There are no new
hypothesis tests. These two calibration families constrain a specific
alternative explanation; they do not exclude nonlinear maps, other features,
training objectives, or larger adaptation samples.
